% =============================================================
%  NONSTOP DESIGN -- schlichter DIN-5008-Geschäftsbrief
%  Vorlage: "Rechnung Kleinstadt Brand.pdf" (Nonstop Performance,
%  Ali Batuhan Arabaci) -- Scan aus den Eingangsrechnungen.
%  Engine: pdfLaTeX
%
%  Alle Maße stammen aus dem Scan. Der Satzspiegel des Originals ist
%  25 mm links / 25 mm rechts (gemessene Kopflinie 159,3 mm breit bei
%  ~0,5 % Scan-Verkleinerung); die gemessenen x-Werte sind darauf
%  normalisiert. Die Schriftgrade sind über Textbreiten kalibriert
%  (Arial/Helvetica sind metrisch kompatibel):
%     Fließtext 10 pt, Zeilenabstand 13,3 pt (4,7 mm)
%     Überschrift 12 pt fett, Zahlungszeile 11 pt
%     Rücksendeangabe + Hinweise 8 pt, Fußzeile 8 pt / 19,3 pt
%
%  Aufbau
%  ------------------------------------------------------------------
%  [x] Logo oben rechts, rechtsbündig zum Satzspiegel (185 mm)
%  [x] Rücksendeangabe (DIN 5008), unterstrichen
%  [x] Anschriftfeld
%  [x] "Rechnungsdatum: ..." rechtsbündig
%  [x] Leistungsdatums-Hinweis (§ 14 Abs. 4 Nr. 6 UStG) -- NICHT im
%      Original; Pflichtangabe, deshalb ergänzt
%  [x] "Rechnung Nr. <Nr.>" fett
%  [x] Anschreiben (Anrede optional, Rechnungstext)
%  [x] Positionstabelle Nr./Bezeichnung/Menge/Einzelbetrag/Gesamt
%      mit durchgehender Linie unter der Kopfzeile
%  [x] Summenblock rechtsbündig: Summe netto / zzgl. X% USt /
%      Rechnungsbetrag (fett); § 19 UStG statt USt-Zeilen
%  [x] Zahlungsbedingung ("sofort fällig" bzw. Zahlungsziel)
%  [x] Dreispaltige Fußzeile: Anschrift | Bankverbindung | Steuerdaten
%  [x] optionaler EPC/Girocode-QR (im Original nicht vorhanden)
%
%  Bewusste Abweichungen vom Scan:
%   - Beträge mit zwei Nachkommastellen (Original: "10.000 €").
%   - Mengen/Beträge stehen auf der ERSTEN Beschreibungszeile
%     (Original: Word-Zellenausrichtung "unten").
%   - "Steuer-Nr.:" / "USt-IdNr.:" statt "Steuer-Nr:" / "Ust.-ID:".
% =============================================================

\documentclass[a4paper,10pt]{article}

% ---- Encoding & Sprache ----------------------------------
\usepackage[T1]{fontenc}
\usepackage[utf8]{inputenc}
\usepackage[ngerman]{babel}
\usepackage{textcomp}

% ---- Satzspiegel -----------------------------------------
% Der Satzspiegel gilt für FOLGESEITEN (Textbeginn 30 mm). Seite 1 trägt
% den Briefkopf als Overlay; ihr Fließtext wird mit \letterheadskip um
% 70 mm nach unten geschoben, sodass die erste Grundlinie exakt auf
% 105,6 mm liegt (Originalmaß). Textunterkante 250 mm, darunter beginnt
% bei 256,4 mm die Fußzeile (footskip = 258,8 mm Grundlinie - 250 mm).
\usepackage[
  a4paper,
  left=25mm,
  right=25mm,
  top=30mm,
  textheight=220mm,
  headheight=0pt,
  headsep=0pt,
  footskip=8.8mm
]{geometry}
\setlength{\topskip}{5.6mm}
% \vspace* setzt eine unsichtbare \hrule der Höhe 0 ab. Die ist die erste
% Box der Seite, bekommt also den \topskip; die Überschrift danach erhält
% zusätzlich einen vollen \baselineskip. Beide werden hier abgezogen,
% damit die erste Grundlinie exakt auf 105,6 mm liegt.
\newcommand{\letterheadskip}{%
  \vspace*{\dimexpr 105.6mm - 30mm - \topskip - \baselineskip\relax}}

% ---- Schrift: Helvetica (metrisch wie das Arial der Vorlage) ----
\usepackage{helvet}
\renewcommand{\familydefault}{\sfdefault}

% ---- Pakete ----------------------------------------------
\usepackage{xcolor}
\usepackage{graphicx}
\usepackage{array}
\usepackage{fp}
\usepackage{etoolbox}
\usepackage{ifthen}
\usepackage{url}
\usepackage{tikz}

\pagestyle{empty}
\setlength{\parindent}{0pt}
\setlength{\parskip}{0pt}
% Das Original ist linksbündig gesetzt (Zeile 1 des Anschreibens endet
% deutlich vor dem rechten Rand) -- kein Blocksatz, keine Trennung.
\raggedright

\newcommand{\bodyfont}{\fontsize{10pt}{13.3pt}\selectfont}
\newcommand{\headfont}{\fontsize{12pt}{14.4pt}\selectfont}
\newcommand{\duefont}{\fontsize{11pt}{13.3pt}\selectfont}
\newcommand{\notefont}{\fontsize{8pt}{9.6pt}\selectfont}
\newcommand{\footfont}{\fontsize{8pt}{19.3pt}\selectfont} % Fußzeile: 6,8 mm Zeilenabstand

% \gapto{<mm>}: vertikaler Abstand so, dass die NÄCHSTE Grundlinie
% genau <mm> unter der aktuellen liegt (Grundlinienabstand abgezogen).
\newcommand{\gapto}[1]{\vspace{\dimexpr #1mm - \baselineskip\relax}}

% (Das gemeinsame _data.tex übergibt Zahlen ohne Füllleerzeichen.)

% =============================================================
%  TABELLENMASSE (aus dem Scan, Satzspiegel 25..185 mm)
% =============================================================
\newlength{\colNr}      \setlength{\colNr}{12.5mm}
\newlength{\colDesc}    \setlength{\colDesc}{65.5mm}
\newlength{\colQty}     \setlength{\colQty}{18mm}
\newlength{\colUnit}    \setlength{\colUnit}{30.3mm}
\newlength{\colTotal}   \setlength{\colTotal}{33.7mm}

% =============================================================
%  BRIEFKOPF-OVERLAY (nur Seite 1)
%  Wird dem ersten Textbaustein vorangestellt; overlay-tikzpictures
%  haben Breite und Höhe 0 und stören den Umbruch deshalb nicht.
% =============================================================
% \atpage{<anchor>}{<x mm>}{<y mm von oben>}{<Inhalt>}
% Mehrzeilige Inhalte werden in \parbox[t] gesetzt: damit ist der
% base-Anker die Grundlinie der ERSTEN Zeile (ein nacktes mehrzeiliges
% tikz-Node verankert sonst auf der letzten Zeile).
\newcommand{\atpage}[4]{%
  \begin{tikzpicture}[remember picture, overlay]
    \node[anchor=#1, inner sep=0pt, outer sep=0pt]
      at ([xshift=#2mm, yshift=-#3mm]current page.north west) {#4};
  \end{tikzpicture}%
}

\newcommand{\letterhead}{%
  % --- Logo oben rechts (Wortmarke als Fallback) -----------
  \ifHasLogo
    \atpage{north east}{185}{30}{\includegraphics[height=16.7mm,keepaspectratio]{logo.png}}%
  \else
    \atpage{base east}{185}{44}{\fontsize{26pt}{28pt}\selectfont\bfseries\shortName}%
  \fi
  % --- Rücksendeangabe (DIN 5008) --------------------------
  \atpage{base west}{25}{48}{%
    \notefont\underline{\senderName\ \textbar\ \senderStreet\ \textbar\ \senderZIP\ \senderCity}}%
  % --- Anschriftfeld ---------------------------------------
  \atpage{base west}{25}{55.8}{%
    \parbox[t]{85mm}{\bodyfont
      \ifdefempty{\customerCompany}{}{\customerCompany\\}%
      \customerName\\
      \customerStreet\\
      \customerZIP\ \customerCity}}%
  % --- Rechnungsdatum --------------------------------------
  \atpage{base east}{185}{78}{\bodyfont Rechnungsdatum:\hspace{3.7mm}\invoiceDate}%
  % --- Leistungsdatum (§ 14 Abs. 4 Nr. 6 UStG) -------------
  \atpage{base east}{185}{82.5}{%
    \notefont
    \ifHasServiceDateNote\serviceDateNote\else
      \ifEnglish Service date equals invoice date.%
      \else Leistungsdatum entspricht dem Rechnungsdatum.\fi
    \fi}%
}

% =============================================================
%  FUSSZEILE (auf jeder Seite)
% =============================================================
\newcommand{\nonstopfooter}{%
  \footfont
  \hspace*{5.55mm}%
  % Spalte 1: Anschrift. Erste Zeile bleibt leer -- im Original beginnt
  % die Anschriftspalte eine Zeile tiefer als Bank und Steuerdaten.
  \parbox[t]{53.65mm}{\mbox{}\\ \senderCompany\\ \senderStreet\\ \senderZIP\ \senderCity}%
  % Spalte 2: Bankverbindung
  \parbox[t]{55.9mm}{%
    \accountBankName\\
    \ifEnglish Account holder:\else Inhaber:\fi\ \accountRCPT\\
    IBAN:~\accountIBAN%
    \ifdefempty{\accountBIC}{}{\\BIC:~\accountBIC}}%
  % Spalte 3: Steuerdaten
  \parbox[t]{44.9mm}{%
    \ifEnglish Tax No.:\else Steuer-Nr.:\fi\ \taxID%
    \ifdefempty{\vatID}{}{\\\ifEnglish VAT ID:\else USt-IdNr.:\fi\ \vatID}%
    \ifHasTaxOffice\\\taxOffice\fi}%
  \hfill
}

\makeatletter
\def\ps@nonstop{%
  \def\@oddhead{}\def\@evenhead{}%
  \def\@oddfoot{\nonstopfooter}\let\@evenfoot\@oddfoot}
\makeatother
\pagestyle{nonstop}

% =============================================================
%  POSITIONEN
%  _invoice.tex ruft \ProjectTitle, \LineDiscount und \Fee auf.
%  Die Zeilen werden in \allFeeRows gesammelt und erst im Dokument
%  gesetzt -- wie im modern-Design inklusive \edef-Expansion, damit
%  spätere \Fee-Aufrufe frühere Zeilen nicht überschreiben.
% =============================================================
\newcounter{itemnum}
\FPset{\invoiceSubtotal}{0}
\newcommand{\allFeeRows}{}

\newcommand{\unitsuffix}[1]{\ifstrempty{#1}{}{\,#1}}

% \FeeNote: optionale Zusatzzeile, vom Go-Layer in das
% Beschreibungsargument eingebettet. \protected\def ist notwendig, weil
% \Fee die Zeile per \edef aufbaut und \newline nicht expandierbar ist.
\protected\def\FeeNote#1{\newline{\notefont #1}}

% \LineDiscount wird vom gemeinsamen Template direkt VOR dem zugehörigen
% \Fee ausgegeben; die Werte werden zwischengespeichert und nach der
% Positionszeile als eigene Zeile gesetzt.
\newcommand{\LDpercent}{}
\newcommand{\LDamount}{}
\def\LDcents{0}
\newcommand{\LineDiscount}[3]{\gdef\LDpercent{#1}\gdef\LDcents{#2}\gdef\LDamount{#3}}

\newcommand{\Fee}[4]{%
  \stepcounter{itemnum}%
  \FPmul{\myLineTot}{#2}{#3}%
  \FPround{\myLineTotR}{\myLineTot}{2}%
  \FPadd{\myNewSub}{\invoiceSubtotal}{\myLineTotR}%
  \FPround{\myNewSubR}{\myNewSub}{2}%
  \xdef\invoiceSubtotal{\myNewSubR}%
  \edef\curRow{%
    \the\value{itemnum} &
    #1 &
    \noexpand\fmtqty{#3}\noexpand\unitsuffix{#4} &
    \noexpand\fmtmoney{#2}\noexpand\,\noexpand\texteuro &
    \noexpand\fmtmoney{\myLineTotR}\noexpand\,\noexpand\texteuro
    \noexpand\\[2.2mm]%
  }%
  \expandafter\gappto\expandafter\allFeeRows\expandafter{\curRow}%
  \ifnum\LDcents=0 \else
    \edef\LDcall{\noexpand\FeeDiscount{\LDpercent}{\LDamount}}\LDcall
    \gdef\LDpercent{}\gdef\LDcents{0}\gdef\LDamount{}%
  \fi
}

% \FeeDiscount: Nachlasszeile unter der Position. Angezeigt wird die
% Minderung; der Rechnungsbetrag kommt aus der Go-Schicht (\NetTotal),
% die den Nachlass bereits abgezogen hat.
\newcommand{\FeeDiscount}[2]{%
  \FPsub{\myNewSub}{\invoiceSubtotal}{#2}%
  \FPround{\myNewSubR}{\myNewSub}{2}%
  \xdef\invoiceSubtotal{\myNewSubR}%
  \gappto\allFeeRows{%
    & \textit{\ifEnglish Discount\else Nachlass\fi\ifstrempty{#1}{}{ #1\,\%}} & & &
    $-$\fmtmoney{#2}\,\texteuro \\[2.2mm]%
  }%
}

\newcommand{\ProjectTitle}[1]{\gdef\theProjectTitle{#1}}
\gdef\theProjectTitle{}

% =============================================================
%  SUMMENZEILE: Beschriftung rechtsbündig in der Einzelbetrag-Spalte,
%  Betrag rechtsbündig in der Gesamt-Spalte.
% =============================================================
\newcommand{\TotalRow}[2]{%
  \makebox[\dimexpr\colNr+\colDesc+\colQty+\colUnit\relax][r]{#1}%
  \makebox[\colTotal][r]{#2}\par}

% =============================================================
\begin{document}

% ################## invoice DATA ##################
\newcommand{\invoiceDate}{26.06.2025} %\today% Datum der rechnungsstellung
\newcommand{\payDate}{11.07.2025} % Datum der Zahlungsfrist
\newcommand{\invoiceReference}{202506-1} % Rechnungsnummer (z.B. 20150122-4)
\newcommand{\invoiceSalutation}{Sehr geehrte Damen und Herren,} % die Anrede
\newcommand{\invoiceText}{für die von mir erbrachte Leistung für das Andreas Heardt Catering am 13.06.2025 erhalten Sie
hiermit die Rechnung.} % Rechnungstext
\newcommand{\invoiceEnclosures}{} % \encl{} einfügen
\newcommand{\invoiceClosing}{Mit freundlichen Grüßen}
% ################## invoice DATA ##################



% ################## Customer DATA ##################
\newcommand{\customerCompany}{} %ggf. Firma
\newcommand{\customerName}{Loser \& Freudenberg GbR} % Name
\newcommand{\customerStreet}{Fischergasse 2} % Straße
\newcommand{\customerZIP}{89073} % Postleitzahl
\newcommand{\customerCity}{Ulm} % Ort
% ################## Customer DATA ##################




% ################## Personal DATA ##################
% START INVOICE DATA
% Einzelunternehmen
%\newcommand{\taxID}{91186 / 08005}
% Catering GbR
\newcommand{\taxID}{91070 / 43733}
% END INVOICE DATA
% START SENDERS DATA
\newcommand{\senderName}{Rico Klatte}
\newcommand{\senderCompany}{RK Holding GmbH\\&und Rico Klatte\\&Catering GbR}
\newcommand{\senderStreet}{Steibisberger Weg 20}
\newcommand{\senderZIP}{88239}
\newcommand{\senderCity}{Wangen im Allgäu}
\newcommand{\senderTelephone}{+49 (0)173 4529776}
\newcommand{\senderMobilephone}{+49 (0)171 4184943}
\newcommand{\senderEmail}{mail@ricoklatte.de}
\newcommand{\senderWeb}{\url{www.ricoklatte.de}}
% END SENDERS DATA
% START ACCOUNT DATA
\newcommand{\accountRCPT}{\senderName}
\newcommand{\accountNumber}{0000000000}
\newcommand{\accountBLZ}{000000000}
\newcommand{\accountBankName}{ING-DiBa}
\newcommand{\accountIBAN}{DE50 5001 0517 5450 3793 40}
\newcommand{\accountBIC}{INGDDEFFXXX}
% END ACCOUNT DATA
% ################## Personal DATA ##################

\ProjectTitle{Kaffee Catering 13.06.2025} %
\Fee{Caffee Crema}{3.5}{1}
\Fee{Latte Macchiato}{4.5}{1}
\Fee{Cappuccino}{4.5}{3}
\Fee{Magnum}{3.5}{4}
\Fee{Doppio}{3.5}{1}
\Fee{Espresso}{2.5}{3}
\Fee{Konfekt}{3}{1}
\Fee{Nogger}{3}{1}
\Fee{Kuchen}{3.5}{2}
\Fee{Stroopwafel}{1.5}{1}
\Fee{Cornetto}{3}{1}
\Fee{Cold Brew}{5}{1}

% Auslagen
%\EBCi{Hotel, 12 Nächte} {2400.00}
%\STExpenses

% Rabatt
%\Discount{Rabatt} {20}

%\Fee{BI Consulting November 2023}{35.00}{12.5}
%\Fee{BI Consulting November 2023}{40.00}{26}
%\Fee{BI Consulting Dezember 2023}{35.00}{3}
%\Fee{BI Consulting Dezember 2023}{40.00}{11}

\bodyfont

% --- Überschrift (trägt das Briefkopf-Overlay) --------------
\letterheadskip
\letterhead
{\headfont\bfseries\ifEnglish Invoice No.\else Rechnung Nr.\fi\ \invoiceReference}\par

% --- Anschreiben --------------------------------------------
\gapto{13.2}
\ifdefempty{\invoiceSalutation}{}{\invoiceSalutation\par\gapto{9.4}}
\ifHasProjectTitle{\bfseries\theProjectTitle\par\gapto{9.4}}\fi
\invoiceText\par

% --- Positionstabelle ----------------------------------------
\gapto{17.0}
{\bfseries
 \makebox[\colNr][l]{Nr.}%
 \makebox[\colDesc][l]{\ifEnglish Description\else Bezeichnung/Leistung\fi}%
 \makebox[\colQty][r]{\ifEnglish Qty\else Menge\fi}%
 \makebox[\colUnit][r]{\ifEnglish Unit price\else Einzelbetrag\fi}%
 \makebox[\colTotal][r]{\ifEnglish Total/\texteuro\else Gesamt/\texteuro\fi}%
}\par
% Linie 0,25 mm UNTER der Grundlinie -- im Original schneidet sie die
% Unterlängen von "Bezeichnung/Leistung".
\nointerlineskip
\vspace{0.25mm}
\rule{\textwidth}{0.4mm}

\vspace{\dimexpr 8.1mm - 0.65mm - \baselineskip\relax}
\setlength{\tabcolsep}{0pt}
\renewcommand{\arraystretch}{1}
\begin{tabular}{@{}p{\colNr}>{\raggedright\arraybackslash}p{\colDesc}%
  >{\raggedleft\arraybackslash}p{\colQty}%
  >{\raggedleft\arraybackslash}p{\colUnit}%
  >{\raggedleft\arraybackslash}p{\colTotal}@{}}
\allFeeRows
\end{tabular}\par

% --- Summenblock ---------------------------------------------
% Kalibriert: erste Summenzeile 12,8 mm unter der letzten Positionszeile
% (Originalmaß). Der Abstand wird ab der UNTERKANTE der Tabelle gemessen
% -- für hohe Boxen greift \lineskip, nicht \baselineskip --, ist also
% unabhängig von der Anzahl der Positionen.
\vspace{8.4mm}
\ifHasAnyVat
  \renewcommand{\VatBreakdownRow}[4]{%
    \TotalRow{\ifEnglish plus #1\% VAT\else zzgl. #1\% Umsatzsteuer\fi}%
             {\fmtmoney{#3}\,\texteuro}}%
  \TotalRow{\ifEnglish Net total\else Summe netto\fi}{\fmtmoney{\NetTotal}\,\texteuro}
  \VatBreakdownRows
  \TotalRow{\bfseries\ifEnglish Invoice total\else Rechnungsbetrag\fi}%
           {\bfseries\fmtmoney{\GrossTotal}\,\texteuro}
\else
  \TotalRow{\ifEnglish Net total\else Summe netto\fi}{\fmtmoney{\NetTotal}\,\texteuro}
  \TotalRow{\bfseries\ifEnglish Invoice total\else Rechnungsbetrag\fi}%
           {\bfseries\fmtmoney{\NetTotal}\,\texteuro}
  \gapto{6.5}
  {\notefont
   \makebox[\textwidth][r]{%
     \ifEnglish
       The amount shown contains no value-added tax pursuant to \S\,19 of the German VAT Act.%
     \else
       Im ausgewiesenen Betrag ist gemäß \S\,19 UStG keine Umsatzsteuer enthalten.%
     \fi}\par}
\fi

% --- Zahlungsbedingung ----------------------------------------
\gapto{16.8}
{\duefont
 \ifPaymentTransfer
   \ifthenelse{\equal{\payDate}{\invoiceDate}\OR\equal{\payDate}{}}%
     {\ifEnglish This invoice is due immediately.\else Die Rechnung ist sofort fällig.\fi}%
     {\ifEnglish Payable without deduction by \payDate.\else Zahlbar ohne Abzug bis zum \payDate.\fi}%
 \fi
 \ifPaymentCashDue
   \ifEnglish Please pay the amount in cash upon delivery.\else Bitte zahlen Sie den Betrag bar bei Übergabe.\fi
 \fi
 \ifPaymentCashPaid
   \ifEnglish Amount received in cash with thanks on \cashPaidDate.\else Betrag dankend in bar erhalten am \cashPaidDate.\fi
 \fi
}\par

% --- Optionale Klauseln ---------------------------------------
\ifboolexpr{bool{ShowWarrantyExcluded} or bool{ShowRetentionOfTitle} or bool{ShowLateFeeWarning}}{%
  \gapto{7.0}
  {\notefont
   \ifShowWarrantyExcluded
     \ifEnglish The sale is made under exclusion of any warranty.%
     \else Der Verkauf erfolgt unter Ausschluss jeglicher Gewährleistung.\fi\par
   \fi
   \ifShowRetentionOfTitle
     \ifEnglish Goods remain property of the seller until full payment.%
     \else Die Ware bleibt bis zur vollständigen Bezahlung Eigentum des Verkäufers.\fi\par
   \fi
   \ifShowLateFeeWarning
     \ifEnglish Late payment incurs statutory interest (5 percentage points above the base rate; 9 percentage points for B2B) pursuant to \S\S\,288, 286 BGB.%
     \else Bei Zahlungsverzug werden gesetzliche Verzugszinsen (5 Prozentpunkte über Basiszinssatz, im B2B 9 Prozentpunkte) gemäß \S\S\,288, 286 BGB erhoben.\fi\par
   \fi}%
}{}

% --- Grußformel (im Original nicht vorhanden, Datenmodell kennt sie) ---
% Die Unterschrift hängt an der Grußformel: der Go-Layer setzt
% \signatureName immer (Default = Absender); ein Name ohne Grußformel
% wäre in diesem Brieflayout nur ein loser Eintrag am Seitenfuß.
% \nobreak hält Grußformel und Unterschrift zusammen: reicht der Platz
% nicht mehr, wandern beide gemeinsam auf die Folgeseite statt die
% Unterschrift allein dort stehen zu lassen.
\ifdefempty{\invoiceClosing}{}{%
  \gapto{9.4}\invoiceClosing\par\nobreak
  \ifdefempty{\signatureName}{}{\gapto{14.1}\nobreak\signatureName\par}%
}

% --- EPC/Girocode-QR (im Original nicht vorhanden) -------------
\ifShowQR
\ifPaymentCashPaid\else
\ifHasQRFile
  \gapto{11.0}
  \includegraphics[width=22mm]{epc-qr.png}\par
  {\notefont\ifEnglish Scan to pay\else Zahlen per Girocode\fi\par}
\fi
\fi
\fi

\end{document}
